\documentclass[11pt,a4paper]{article}
\usepackage[UTF8]{ctex}
\usepackage{geometry}
\geometry{left=2.2cm,right=2.2cm,top=2.4cm,bottom=2.4cm}
\usepackage{graphicx}
\usepackage{booktabs}
\usepackage{siunitx}
\usepackage{hyperref}
\usepackage{float}
\usepackage{amsmath}
\usepackage{enumitem}
\usepackage{longtable}
\usepackage{xcolor}

\title{三光子时空稀疏采样快速高质量重建：Stage-3 系统研究报告}
\author{项目：three-photon-spatiotemporal-sparse-sampling-fast-high-quality-reconstruction}
\date{\today}

\begin{document}
\maketitle

\begin{abstract}
本报告对 threePM 课题在 Stage-3 阶段的研究工作进行系统整理，覆盖研究背景、研究目的、总体实验设计、具体实验方法、结果展示与分析、总结与讨论。实验主线聚焦于 2.5D 时空重建方案（以相邻时间窗帧作为输入，重建中心帧），并与 2D baseline 及不同窗口设置进行比较。基于当前 Stage-3 全量评测结果，模型在结构指标上相较噪声输入有提升，但在核心目标“相对既有基线持续提升”上未达到里程碑通过标准，外部审查结论为 Not Pass。报告最后给出下一阶段最小可执行改进路径。
\end{abstract}

\section{研究背景}
三光子显微成像在深层组织观察中具有重要价值，但在实际采集中受限于光毒性、采集时间和信噪比，常采用稀疏采样策略。稀疏采样虽提高采集效率，但会导致图像质量下降，进而影响后续定量分析。该课题目标是在给定稀疏观测条件下，恢复高质量目标图像，兼顾重建质量与时空一致性。

现有 2D 重建方法通常按帧独立处理，忽略时序邻域信息；而 3D 全体积方法对算力与数据组织要求更高。2.5D 方案在工程上具有折中优势：利用局部时序上下文提升重建质量，同时保持网络与训练流程可控。

\section{研究目的}
本阶段目标如下：
\begin{enumerate}[leftmargin=2em]
    \item 在完整评测集上验证 2.5D 方案（$w=3$）在长预算训练（180 epochs）下的稳定性；
    \item 对比 Stage-2 pilot 结果，分析窗口大小与训练预算对重建效果的影响；
    \item 形成可追踪的定量与定性证据，为后续论文写作与实验决策提供依据；
    \item 接入外部审查（GLM-5）并据此给出下一阶段执行清单。
\end{enumerate}

\section{总体实验设计}
\subsection{数据与任务定义}
\begin{itemize}[leftmargin=2em]
    \item 数据集：CIDC25\_paired\_A1F0（本阶段使用 full split 进行评测汇总）。
    \item 输入输出：低质量（noisy/sparse）图像到高质量（dense/clean）图像重建。
    \item 任务形态：\texttt{spatial\_2p5d}，输入时间窗口为 $[t-1,t,t+1]$（$w=3$），输出中心帧 $t$。
\end{itemize}

\subsection{对比设置}
为保证横向可解释性，本报告纳入以下结果：
\begin{itemize}[leftmargin=2em]
    \item Baseline-2D (pilot)
    \item 2.5D $w=3$ (pilot)
    \item 2.5D $w=5$ (pilot)
    \item 2.5D $w=3$ (Stage-3 full, 180 epochs)
\end{itemize}

\subsection{评估指标}
主要指标包括 PSNR、SSIM、MS-SSIM、MAE、LPIPS、$\Delta$PSNR、$\Delta$SSIM、sSNR、FwPSNR。其中特别关注：
\begin{itemize}[leftmargin=2em]
    \item 绝对重建质量：PSNR/SSIM/MS-SSIM/LPIPS；
    \item 相对输入收益：$\Delta$PSNR 与 $\Delta$SSIM；
    \item 稳定性分布：样本级指标直方图与分位样例。
\end{itemize}

\section{具体实验方法}
\subsection{模型与训练}
采用 Swin-UNet 2D 骨干，在 2.5D 任务中将时间窗帧堆叠为输入通道。训练时采用混合损失（包含像素项、结构项及感知项），并保留 DINO 分支。Stage-3 运行配置为：
\begin{itemize}[leftmargin=2em]
    \item 运行名：\texttt{stage3\_2p5d\_w3\_e180\_full}
    \item epoch 数：180
    \item 输出评测目录：\texttt{eval\_results/stage3\_2p5d\_w3\_e180\_full\_eval/eval}
\end{itemize}

\subsection{可视化生成策略}
本报告图表由脚本自动生成（\texttt{scripts/generate\_stage3\_figures.py}），输出到：
\texttt{projects/three-photon-spatiotemporal-sparse-sampling-fast-high-quality-reconstruction/figures\_stage3/}。

图表包括：
\begin{itemize}[leftmargin=2em]
    \item 跨阶段关键指标对比图；
    \item Stage-3 样本级分布图（PSNR/SSIM/$\Delta$PSNR）；
    \item Stage-3 定性重建样例图（低分位/中位/高分位）。
\end{itemize}

\section{结果展示与分析}
\subsection{核心定量结果（Stage-3 full）}
根据汇总文件 \texttt{summary\_stage3\_2p5d\_w3\_e180\_full.csv}：
\begin{longtable}{lS[table-format=2.3]}
\toprule
指标 & {均值} \\
\midrule
PSNR & 25.125 \\
SSIM & 0.639 \\
MS-SSIM & 0.594 \\
MAE & 0.0505 \\
LPIPS & 0.630 \\
$\Delta$PSNR & -0.535 \\
$\Delta$SSIM & 0.245 \\
sSNR & 1.359 \\
FwPSNR & 25.105 \\
\bottomrule
\end{longtable}

解读：模型对结构相似性（$\Delta$SSIM）有提升，但像素级峰值信噪比相对输入未提升（$\Delta$PSNR 为负），且 LPIPS 偏高，说明感知质量仍有改进空间。

\subsection{跨阶段对比}
\begin{figure}[H]
    \centering
    \includegraphics[width=0.95\textwidth]{figures_stage3/fig_metric_comparison.png}
    \caption{跨阶段指标对比（Baseline-2D、2.5D pilot、Stage-3 full）}
    \label{fig:metric_compare}
\end{figure}

从图 \ref{fig:metric_compare} 可见，2.5D 在 pilot 条件下可获得一定增益，但 Stage-3 full 结果未复现到期望水平，提示存在数据分布差异、训练策略或损失权重平衡问题。

\subsection{样本分布与稳定性}
\begin{figure}[H]
    \centering
    \includegraphics[width=0.98\textwidth]{figures_stage3/fig_stage3_distributions.png}
    \caption{Stage-3 样本级分布（PSNR、SSIM、$\Delta$PSNR）}
    \label{fig:dist}
\end{figure}

分布图显示结果方差较大，存在一批样本重建收益不足；尤其是 $\Delta$PSNR 分布含较多负值区间，说明模型在部分场景下未稳定优于输入。

\subsection{定性结果}
\begin{figure}[H]
    \centering
    \includegraphics[width=0.98\textwidth]{figures_stage3/fig_stage3_qualitative_examples.png}
    \caption{Stage-3 定性样例（低分位/中位/高分位）}
    \label{fig:qual}
\end{figure}

定性图可见：高分位样本中结构恢复较好，而低分位样本仍存在细节偏差与对比度恢复不足，和定量分布结论一致。

\subsection{外部审查结论}
根据 \texttt{STAGE3\_REVIEW\_GLM5.md}，外部审查给出 \textbf{Not Pass}，核心依据包括：
\begin{itemize}[leftmargin=2em]
    \item 绝对质量指标未达到预期（尤其 SSIM、LPIPS）；
    \item 相对改进不一致（$\Delta$PSNR 为负）；
    \item 指标间存在不一致，需进一步排查训练与损失权重配置。
\end{itemize}

\section{总结与讨论}
\subsection{阶段结论}
Stage-3 已完成完整训练与评测收敛，并形成了可复现实验产物与图表证据；但就里程碑判定而言，\textbf{当前阶段不通过}。因此，当前产物适合作为“中期研究记录”，尚不宜冻结为最终论文主结果。

\subsection{问题定位}
主要问题集中在以下三点：
\begin{enumerate}[leftmargin=2em]
    \item 感知质量偏弱（LPIPS 高）；
    \item 样本级收益不稳定（$\Delta$PSNR 分布跨越较大且含负区间）；
    \item Stage-2 到 Stage-3 的泛化迁移不足（pilot 优势未稳定延续到 full split）。
\end{enumerate}

\subsection{下一步建议（最小可执行）}
与 \texttt{STAGE3\_NEXT\_ACTIONS.md} 一致，建议按以下顺序推进：
\begin{enumerate}[leftmargin=2em]
    \item 复核训练/验证曲线，确认是否过拟合或欠拟合；
    \item 调整损失项权重（重点感知分支与结构分支平衡）；
    \item 在同一 split 上做严格 paired 对比（Stage-2/2D vs Stage-3/2.5D）；
    \item 增加一个窗口或融合策略消融；
    \item 更新报告并再次提交外部审查。
\end{enumerate}

\section*{附录：文件索引}
\begin{itemize}[leftmargin=2em]
    \item 阶段报告：\texttt{STAGE3\_REPORT.md}
    \item 审查意见：\texttt{STAGE3\_REVIEW\_GLM5.md}
    \item 后续动作：\texttt{STAGE3\_NEXT\_ACTIONS.md}
    \item 本文图表目录：\texttt{figures\_stage3/}
\end{itemize}

\end{document}
